\subsection{Absolute retracts from hyperspace orbits}\label{subsec:ar-hyperspace-orbits}
\paraid{p-bb0d31d746c9}
Equation \eqref{eq:dual-equivariance} converts a change of coordinates into an
isometry of a fixed Banach space.
The global approximation will use compact subsets of such a space
modulo this action.
We apply \Cref{thm:equivariant-hyperspace,thm:orbit-retract} to prove that the resulting hyperspace orbit is
an absolute retract.

\begin{theorem}\label{lem:hyperspace-orbit-ar}
\paraid{p-4a30396f26c4}
Let
$(\BanachAmbientSpace ,\norm{\cdot}_\BanachAmbientSpace )$
be a separable Banach space,
and let a compact metrizable group
$\ActingGroup $
act continuously on
$\BanachAmbientSpace $
by linear isometries.
Equip $\BanachAmbientSpace$ with its induced metric
$\BanachMetric\colon\BanachAmbientSpace\times\BanachAmbientSpace\to[0,\infty)$.
For every $\FirstVector,\SecondVector\in\BanachAmbientSpace$, it satisfies
$\BanachMetric(\FirstVector,\SecondVector)=\norm{\FirstVector-\SecondVector}_\BanachAmbientSpace$.
Then the orbit space
$\CompactHyperspace(\BanachAmbientSpace )/\ActingGroup $
is a separable metrizable AR.
For nonempty compact subsets
$\FirstCompactSet,\SecondCompactSet\subset\BanachAmbientSpace$,
the formula
\begin{equation}\label{eq:orbit-metric}
  \OrbitMetric([\FirstCompactSet ],[\SecondCompactSet ])=\min_{\GroupElement \in \ActingGroup }\HausdorffDistance{\BanachMetric }(\FirstCompactSet ,\GroupElement \SecondCompactSet )
\end{equation}
defines a metric that induces the quotient topology.
\end{theorem}

\begin{proof}
\textbf{Step 1. The absolute retract property.}
\paraid{p-55569123ecd5}
The induced action on
$\CompactHyperspace(\BanachAmbientSpace )$
is continuous.
To see this,
assume that
$\GroupElement _\SequenceIndex \to \GroupElement $
and
$\HausdorffDistance{\BanachMetric }(\FirstCompactSet _\SequenceIndex ,\FirstCompactSet )\to0$.
Then
\[
 \HausdorffDistance{\BanachMetric }(\GroupElement _\SequenceIndex \FirstCompactSet _\SequenceIndex ,\GroupElement \FirstCompactSet )
 \leq\HausdorffDistance{\BanachMetric }(\FirstCompactSet _\SequenceIndex ,\FirstCompactSet )+\sup_{\CompactSetPoint \in \FirstCompactSet }\norm{\GroupElement _\SequenceIndex \CompactSetPoint -\GroupElement \CompactSetPoint }_\BanachAmbientSpace \to0.
\]
The last supremum tends to
$0$
because the action is continuous and
$\FirstCompactSet $
is compact.
Each group element acts isometrically on the hyperspace
$\CompactHyperspace(\BanachAmbientSpace)$.
By \Cref{thm:compact-hyperspace-topology},
the Hausdorff distance induces the Vietoris topology on
$\CompactHyperspace(\BanachAmbientSpace)$.
\Cref{thm:equivariant-hyperspace} therefore applies because
$\BanachAmbientSpace $
is connected and locally continuum-connected.
Thus
$\CompactHyperspace(\BanachAmbientSpace )$
is a
$\ActingGroup $-AR.
The compact metrizable group
$\ActingGroup $
has a countable basis,
so \Cref{thm:orbit-retract} implies that
$\CompactHyperspace(\BanachAmbientSpace )/\ActingGroup $
is an AR in the metrizable category.

\textbf{Step 2. The orbit metric and quotient topology.}
\paraid{p-154b529a9133}
Let
$\FirstCompactSet,\SecondCompactSet\subset\BanachAmbientSpace$
be nonempty compact subsets.
The minimum in \eqref{eq:orbit-metric} exists because
$\GroupElement \mapsto\HausdorffDistance{\BanachMetric }(\FirstCompactSet ,\GroupElement \SecondCompactSet )$
is continuous on the compact group
$\ActingGroup$.

\paraid{p-fe18da324b94}
The Hausdorff distance on
$\CompactHyperspace(\BanachAmbientSpace)$
is
$\ActingGroup$-invariant,
and every orbit is compact and hence closed.
The argument in
\cite[proof of Theorem~4.3.4, p.~319]{Palais1961}
(see also \cite[p.~4, (2.1) and the preceding paragraph]{Antonyan2020Euclidean})
shows that \eqref{eq:orbit-metric} defines a metric inducing the quotient topology.
Let
$\OrbitProjection \colon\CompactHyperspace(\BanachAmbientSpace )\to\CompactHyperspace(\BanachAmbientSpace )/\ActingGroup $
be the orbit map.

\paraid{p-1e603606de4c}
Since
$\BanachAmbientSpace$
is separable,
\Cref{thm:compact-hyperspace-topology} implies that
$\CompactHyperspace(\BanachAmbientSpace)$
is separable.
Its continuous image under
$\OrbitProjection$
is therefore separable as well.
\end{proof}

\begin{lemma}\label{lem:hyperspace-orbit-realization}
\paraid{p-0459ff3d5b20}
Let
$(\BanachAmbientSpace ,\norm{\cdot}_\BanachAmbientSpace )$
be a separable Banach space,
and let a compact metrizable group
$\ActingGroup $
act continuously on
$\BanachAmbientSpace $
by linear isometries.
Equip $\BanachAmbientSpace$ with its induced metric
$\BanachMetric\colon\BanachAmbientSpace\times\BanachAmbientSpace\to[0,\infty)$.
For every $\FirstVector,\SecondVector\in\BanachAmbientSpace$, it satisfies
$\BanachMetric(\FirstVector,\SecondVector)=\norm{\FirstVector-\SecondVector}_\BanachAmbientSpace$.
Equip
$\CompactHyperspace(\BanachAmbientSpace)/\ActingGroup$
with the metric in \eqref{eq:orbit-metric}.
Then the map
\[
 \RealizationMap\colon\CompactHyperspace(\BanachAmbientSpace )/\ActingGroup \to\GHSpace
\]
defined by
\[
 \RealizationMap([\FirstCompactSet ])=\MetricSpace{\FirstCompactSet }{\RestrictedMetric{\BanachMetric }{\FirstCompactSet }},
\]
is well-defined and
$1$-Lipschitz.
\end{lemma}

\begin{proof}
\paraid{p-0ad0a6a4a958}
For each
$\GroupElement \in \ActingGroup $,
the map
$\SecondCompactSet \to \GroupElement \SecondCompactSet $
obtained by restricting
$\GroupElement$
is an isometry.
Consequently
$\RealizationMap$
is independent of the representative of an orbit.
For every
$\GroupElement\in\ActingGroup$,
we also have
\[
 \begin{aligned}
 \GHDistance(\RealizationMap([\FirstCompactSet ]),\RealizationMap([\SecondCompactSet ]))
 &=\GHDistance(\MetricSpace{\FirstCompactSet }{\BanachMetric |_{\FirstCompactSet \times \FirstCompactSet }},
         \MetricSpace{\GroupElement \SecondCompactSet }{\BanachMetric |_{\GroupElement \SecondCompactSet \times \GroupElement \SecondCompactSet }})\\
 &\leq\HausdorffDistance{\BanachMetric }(\FirstCompactSet ,\GroupElement \SecondCompactSet ).
 \end{aligned}
\]
Taking the minimum over
$\GroupElement\in\ActingGroup$
proves that
$\RealizationMap$
is
$1$-Lipschitz.
\end{proof}

\begin{remark}\label{rem:peano-hyperspace}
\paraid{p-15a612ab8c6c}
A classical hyperspace theorem also relates compact subsets to the Hilbert cube.
Let
$(X,d)$
be a nondegenerate Peano continuum,
that is,
a compact connected locally connected metric space with at least two points.
Curtis and Schori \cite[Theorem~3.2]{CurtisSchori1978}
prove that
\[
 \bigl(\CompactHyperspace(X),\HausdorffDistance{d}\bigr)
 \text{ is homeomorphic to }\HilbertCube=[0,1]^{\NonnegativeIntegers}.
\]
This hyperspace consists of all nonempty compact subsets of the fixed space
$X$.
The space in \Cref{lem:hyperspace-orbit-ar}
instead consists of compact subsets of a Banach space modulo a compact group.
\end{remark}
